% state/admissibility.tex

\chapter{Admissibility States}

\label{chap:admissibility-states}

Admissibility is not an additional geometric state layer. It is the
purpose-qualified evaluation that follows the dependency chain developed in
the preceding chapters. It asks whether derived quantities remain inside an
applicable operational envelope; it neither changes the running Alignment nor
turns a computed result into a decision.

\section{Admissibility State}

\begin{definition}[Admissibility state]
\label{def:admissibility-state}
An admissibility state is a runtime evaluation state describing whether
operational quantities remain within allowed envelopes.
\end{definition}

Admissibility states may include bounds on effective acceleration, jerk,
cant deficiency, roll evolution, velocity, and realization quality.

These quantities are derived from runtime railway state and dynamic
evaluation; they are not persisted as alignment identity.

\section{Operational Envelopes}

\begin{definition}[Operational envelope state]
\label{def:operational-envelope-state}
An operational envelope state describes admissible regions in runtime
state space.
\end{definition}

Operational envelopes are typically defined by standards,
infrastructure constraints, and operational policy. In AIM they act as
constraints on derived runtime quantities rather than as replacements
for intrinsic geometry or railway state.

% ============================================================
\section{Layer Position}
% ============================================================

Admissibility is evaluated after pose2/pose3 reconstruction, metric
realization, frame construction, and speed-qualified quantity derivation.
Rule, version, purpose, applicable object and state, validity time, and
authority must be stated. The assessment remains downstream of railway state
and does not redefine it.

Vehicle-level admissibility is evaluated relative to railway state and
must not be conflated with the railway-state layer itself.

\section{Connection to AIM}

\begin{summary}
Admissibility states connect runtime dynamics, engineering rules,
realization-aware behaviour, and optimization constraints while
preserving the separation between persistent alignment information and
derived operational interpretation.
\end{summary}
